\chapter{Design choices, schematics, simulations and physical implementation}

\begin{figure}[H]
    \center
    \includesvg[height=0.81\textheight]{media/Design}
    \caption{Hierarchical schematic of the implemented mixer}
    \label{fig:sch_root}
\end{figure}

The chosen architecture for the project is the split-mixer architecture, with
only channel modules and a master module. Since the project also implements the
functionality of an audio interface, groups and auxiliary inputs and outputs
have been omitted.

Figure \ref{fig:sch_root} shows the modules of the implemented mixer and how
they will be connected together.
The three channel inputs and the output are connected with consumer line lebel
unbalanced audio RCA jacks. All the modules are implemented on the same PCB
and, as such, are connected together with tracks.

The following design considerations will apply to all the modules described in
the following sections:
\begin{itemize}
    \item All opamp ICs presented are part of the LMV3xx family,
        either the LMV358 or  LMV324 which are dual opamps and quad opamps
        respectively.\cite{lmv3xx}
    \item The whole board will be powered from the USB connection, a maximum
        voltage of 5V and maximum current of 500mA. An LDO regulator is used to
        power all the circuits at 3.3V.
    \item The internal level used by the mixer was chosen to be -16dBV, lower
        than the line level. This is necessary to avoid clipping.
    \item The 3.3V supply of the opamps converts to an ideal maximum signal
        level of 1.17V\sub{RMS} or +1.36dBV. This translates into a headroom of
        17.36dB. All the stages in the modules were calibrated for a maximum
        gain of 10dB, giving a maximum implemented signal level of -6dBV, a
        level that is theoretically low enough to minimise the risk of
        clipping.
\end{itemize}

\section{The Channel Module}

\begin{figure}[H]
    \center
    \includesvg[width=\textwidth]{media/Design-PREAMP-EQ_L1}
    \caption{Hierarchical schematic of half of the stereo channel module}
    \label{fig:sch_channel}
\end{figure}

The implemented channel module seen in figure \ref{fig:sch_channel} consists of
an input stage and an EQ stage, interconnected with the stereo audio codec.

The input stage begins with a DC blocking capacitor (C19 in the figure)
connected to a voltage divider (R9 and R10) further connected to the input of
an inverting amplifier (U2A). The DC blocking capacitor and the voltage divider
add a DC component to the input signal, while the inverting amplifier presents
an adequate impedance for a line in connection and brings the signal level down
to.
The inverse phase shifting
introduced by the first opamp is then corrected by the second opamp (U2B) in
the stage, which also controls the gain of the signal through the potentiometer
(RV1). The output of this stage is then sent to the audio codec section.

The EQ stage receives its signals from the output of the audio codec.
It consists of three parallel integration and derivation opamp configurations
(U3A, B and C), each of them acting as a band pass filter, with bands 20Hz to
200Hz, 200Hz-2kHz, 2kHz-20kHz. The fourth opamp (U3C) is in an inverting summer
configuration, adding together the outputs of the three filters. Fader RV5 acts
as a volume control for the channel, before the mixing stage.

\begin{figure}[H]
    \center
    \includegraphics[width=\textwidth]{media/sim_preamp-ac}
    \caption{Multi-step AC simulation of the input stage}
    \label{fig:sim_preamp-ac}
\end{figure}

Figure \ref{fig:sim_preamp-ac} shows the frequency response of the input stage
for different wiper positions on the potentiometer.
The 0dB reference in this graph is line level. The different plots have no
distinct curvature, meaning the audio signals pass through the circuit
unaltered except for amplitude, as intended.

\begin{figure}[H]
    \center
    \includegraphics[width=\textwidth]{media/sim_preamp-in}
    \caption{Simulation of the input impedance of the input stage}
    \label{fig:sim_preamp-in}
\end{figure}

\begin{figure}[H]
    \center
    \includegraphics[width=\textwidth]{media/sim_preamp-out}
    \caption{Simulation of the output impedance of the input stage}
    \label{fig:sim_preamp-out}
\end{figure}

Figures \ref{fig:sim_preamp-in} and \ref{fig:sim_preamp-out} show the
input and output impedances of the stage varying with the frequency.
Each curve corresponds to a different setting of the potentiometer.
Both graphs suggest that the impedances are within the specifications.

\begin{figure}[H]
    \center
    \includegraphics[width=\textwidth]{media/sim_preamp-noise}
    \caption{Simulation of the noise spectrum in the preamp stage}
    \label{fig:sim_preamp-noise}
\end{figure}

Using the available noise analysis tools in SIMetrix, the power spectrum can
be obtained, seen in figure \ref{fig:sim_preamp-noise} and the total noise of
this stage can be computed and will be -103.46dBV.

\begin{figure}[H]
    \center
    \includegraphics[width=\textwidth]{media/sim_eq-ac-low}
    \caption{Low frequencies boost/cut}
    \label{fig:sim_eq-low}
\end{figure}

\begin{figure}[H]
    \center
    \includegraphics[width=\textwidth]{media/sim_eq-ac-mid}
    \caption{Mid frequencies boost/cut}
    \label{fig:sim_eq-mid}
\end{figure}

\begin{figure}[H]
    \center
    \includegraphics[width=\textwidth]{media/sim_eq-ac-high}
    \caption{High frequencies boost/cut}
    \label{fig:sim_eq-high}
\end{figure}

Figures \ref{fig:sim_eq-low}, \ref{fig:sim_eq-mid} and \ref{fig:sim_eq-high}
show the frequency response of the EQ section for different frequency bands and
different potentiometer settings. Since the integration and derivation opamp
configuration is, in essence, a high-pass filter followed by a low-pass filter,
it's only natural that towards the ends of the spectrum the amplitude will
always be lower. The location of the poles and zeroes of the transfer functions
also shift with the potentiometer settings.

\begin{figure}[H]
    \center
    \includegraphics[width=\textwidth]{media/sim_eq-in}
    \caption{Input impedance of the EQ stage}
    \label{fig:sim_eq-in}
\end{figure}

\begin{figure}[H]
    \center
    \includegraphics[width=\textwidth]{media/sim_eq-out}
    \caption{Output impedance of the EQ stage}
    \label{fig:sim_eq-out}
\end{figure}

Figures \ref{fig:sim_eq-in} and \ref{fig:sim_eq-out} show the input and output
impedances of the EQ stage, varying with frequency. The input impedance is
lower than typical for line level inputs, but not such that it will
categorically cause overloading of the codec's output. The output impedance on
on the other hand (measured at the opamp output, before the fader) is more than
adequate to drive the following stage.

\begin{figure}[H]
    \center
    \includegraphics[width=\textwidth]{media/sim_eq-noise}
    \caption{Simulation of the noise spectrum in the EQ stage}
    \label{fig:sim_eq-noise}
\end{figure}

Using the available noise analysis tools in SIMetrix, the power spectrum can
be obtained, seen in figure \ref{fig:sim_eq-noise} and the total noise of
this stage can be computed and will be -105.69dBV.


\section{The Master Module}

\begin{figure}[H]
    \center
    \includesvg[width=\textwidth]{media/Design-MASTER_L}
    \caption{Hierarchical schematic of half of the master module}
    \label{fig:sch_master}
\end{figure}

The master module presented in \ref{fig:sch_master}, similarly to the channel
modules, begins with DC blocking capacitors and resistive voltage dividers which
re-center the received signal on the necessary DC component. Then the inverting
summer (opamp U16A in the figure) adds together the three channels, and sends
the resulting signal through the fader acting as a volume control (RV16) and
further on, to a final opamp stage (U16B) which presents and adequate impedance
for a line out connection.

\begin{figure}[H]
    \center
    \includegraphics[width=\textwidth]{media/sim_master-ac}
    \caption{AC simulation of the master stage}
    \label{fig:sim_master-ac}
\end{figure}

As with the preamplifier stage, the frequency response of the master module
needs to be as flat as possible, since it is only meant to sum the signals
coming from the three channels. Figure \ref{fig:sim_master-ac} shows the
frequency response of the stage for a single active input. If all inputs were
active, the resulting curve would show a gain of 3 or roughly 10dB
(the sum of the three AC voltage sources, one connected to each input, being
three times higher than the signal coming through a single channel).
In practical terms, this would be a reflection of three signals that occupy the
whole audio spectrum, raising the chance of clipping. This is slightly
unrealistic in a studio scenario:
Ideally, in a studio, each channel (representing different instruments or
voices) will occupy a different frequency band, the result being a signal that
occupies the whole spectrum. A channel outputting a signal that occupies the
whole spectrum would mean a song, or something similar, is playing
through that channel, something that is more likely to happen in a live
scenario, such as PA or radio broadcast, where accompanying music would be
played.

\begin{figure}[H]
    \center
    \includegraphics[width=\textwidth]{media/sim_master-in}
    \caption{Simulation of the input impedance of the master stage}
    \label{fig:sim_master-in}
\end{figure}

\begin{figure}[H]
    \center
    \includegraphics[width=\textwidth]{media/sim_master-out}
    \caption{Simulation of the output impedance of the master stage}
    \label{fig:sim_master-out}
\end{figure}

Figures \ref{fig:sim_master-in} and \ref{fig:sim_master-out} show the input and
output impedances of this stange. As in the case of the channel module, the
output impedance is within the specifications of line level output, and the
input impedance is more than adequate to be driven by the EQ stage that
precedes it.

\begin{figure}[H]
    \center
    \includegraphics[width=\textwidth]{media/sim_master-noise}
    \caption{Simulation of the noise spectrum in the master module}
    \label{fig:sim_master-noise}
\end{figure}

Using the available noise analysis tools in SIMetrix, the power spectrum can
be obtained, seen in figure \ref{fig:sim_master-noise} and the total noise of
this stage can be computed and will be -106.61dBV.


\section{Audio Codec Section}

\begin{figure}[H]
    \center
    \includesvg[width=\textwidth]{media/Design-PCM1}
    \caption{Hierarchical schematic of the audio codec circuitry}
    \label{fig:sch_codec}
\end{figure}

This section is realised using the PCM2903C integrated circuit.
The PCM2903C is a single-chip USB audio codec designed for various audio
applications, including audio recording and playback devices.
It integrates USB transceiver functionality, stereo ADC, stereo DAC, and an
audio interface into a single device.

The USB audio interface is compliant with USB 2.0 full speed technology (a
maximum transfer rate of 12Mbps) and supports USB HID (Human Interface Device)
for audio controls.

The ADC and DAC both support 16-bit resolution and sampling rates of up to
48kHz. In terms of audio performance, the typical SNR values are 89dB for the
ADC and 96dB for the DAC and the dynamic ranges are 89dB and 93dB respectively.
Both the ADC and the DAC conform with
consumer line level signals, and present the necessary impedances.

In terms of power management, it is powered from single supply of 3.3V and
draws a current of 54mA during operation.
\cite{pcm}

\begin{figure}[H]
    \center
    \includegraphics[width=\textwidth]{media/pcm-startup}
    \caption{Startup of the PCM2903C}
    \label{fig:pcm-startup}
\end{figure}

For the implementation, seen in figure \ref{fig:sch_codec} the power input pins
are all connected with decoupling
capacitors to filter low frequency power supply noise. The integrated circuit
features a crystal oscillator and requires an external crystal resonator
(Y1 in the schematic) in order to function. The data transmission is
handled by a differential pair, connected at pins 1 and 2 (D+ and D-) of the
device. Both pins are connected to decoupling capacitors and the D+ also
features a pullup resistor for correct functioning. Pin SEL1 triggers the
integrated circuit's reset procedure, preparing it for functioning, as seen
in figure \ref{fig:pcm-startup}, as such it is connected to an RC network to
delay the procedure until after the power is supplied to the circuit.
The analog audio pins all feature DC blocking capacitors to prevent any DC
components or noise from passing through to the ADC or out from the DAC.

\section{USB Hub Section}

\begin{figure}[H]
    \center
    \includesvg[width=0.8\textwidth]{media/Design-USB HUB}
    \caption{Hierarchical schematic of the USB hub circuitry}
    \label{fig:sch_usbhub}
\end{figure}

The TUSB2046 is a 4-port USB 1.1 hub controller that allows for the expansion
of a single USB port into four downstream ports. It is designed for embedded
systems and peripheral devices requiring multiple USB connections. The device
complies with the USB 1.1 standard and supports both full-speed (12 Mbps).

In terms of power management, the integrated circuit is powered from a single
supply of 3.3V and draws a maximum current of 40mA.
\cite{tusb}
All power pins are similarly connected with decoupling capacitors to filter
power supply noise, as seen in figure \ref{fig:sch_usbhub}.
This integrated circuit also features a crystal oscillator
that requires an external crystal resonator to function properly (Y4 in the
schematic). Since only three channels are implemented, only three of the
downstream data pins of this device are used, the last differential pair is
connected to ground with two resistors (pins 23 and 24). All the downstream
data pins are otherwise connected simliarly to the PCM2903C data pins, with
decoupling capacitors and pull down resistors. The upstream data pins again
follow the same pattern and are further connected to the mini USB type B port.
Pin 4 of this device acts as enable pin and is also connected to an RC network
that controls delays the startup of circuit.
In terms of transfer rate, since each codec produces 16-bit stereo samples and
48000 samples per second, the total necessary bandwidth is 9.2Mbps for
bidirectional communication, less than the available 12Mbps offered by the
integrated circuits. Furthermore, the total delay introduced by the digital
communication on the board is 4ms, 2 upstream transfers and 2 downstream
transfers, each of needing 1ms to be accomplished, as defined by the USB full
speed standard\cite{usb20}.

For power, the MCP1703A low dropout voltage regulator was used. It supplies the
whole board with 3.3V obtained from the 5V bus pin of the USB port. It supplies
a maximum current of 250mA, more than the 202mA necessary for the operation of
the codecs and the usb hub circuits, with enough room to power the analog
circuitry as well.

\section{PCB Layout}

\vfill
\begin{figure}[H]
    \center
    \includegraphics[width=\textwidth]{media/PCB_front}
    \caption{Front view of the PCB}
    \label{fig:pcb_front}
\end{figure}
\vfill
\begin{figure}[H]
    \center
    \includegraphics[width=\textwidth]{media/PCB_back}
    \caption{Back view of the PCB}
    \label{fig:pcb_back}
\end{figure}

The PCB reflects the split-mixer architecture, the three channels are placed to
the left of the master module, the total size is 200mm by 146.5mm.
It features ground planes on both copper layers,
screening most of the tracks from crosstalk.

The front layer features all the digital components on the board while to back
layer features most of the analog components, effectively isolating the two
return paths of the signals from each other, to minimise the noise.

The three cannel modules are also copies of each
other, and all operational amplifiers are process only a single side of the
stereo audio signals, minimising crosstalk between the stereo channels and
between the two sides of each individual channel.

Furthermore, the tracks are 1oz copper (35 microns thick) and their widths are
0.3mm for tracks carrying signals and 0.4mm for power supply tracks. This
results in a resistance of
$1.6\frac{\Omega}{\text{m}}$ for the signal tracks and
$1.2\frac{\Omega}{\text{m}}$ for the power supply tracks.
Given the dimensions of the board, the maximum length of any straight track
can't be more than 200mm, and the resistance of that track will be at most
0.32\textOmega for signal tracks. Given the fact that most input impedances
were designed to be in the magnitude of tens of kiloOhms, voltage mismatch
between stages as a result of track resistance will be minimal.

For the differential pairs between the digital components, care was taken such
that each pair is screened from others and that they cross as few of the
analog signal tracks as possible.

All passive components were placed as close as possible to the integrated
circuits on both sides of the board, the exception being the potentiometers and
the faders.

As of the 27th of June, the board was not completely assembled. Out of the 410
components, only around 150 were placed on the board, covering one channel
module with the associated codec section and the usb hub section. Figures
\ref{fig:statfront} and \ref{fig:statback} show the partial assembly of the
board.

\begin{figure}[h]
    \center
    \includegraphics[width=\textwidth]{media/pcblicenta2}
    \caption{Assembly of the front side of the PCB}
    \label{fig:statfront}
\end{figure}

\begin{figure}[h]
    \center
    \includegraphics[width=\textwidth]{media/pcblicenta1}
    \caption{Assembly of the back side of the PCB}
    \label{fig:statback}
\end{figure}



